% GENERATED FILE. Edit canonical lessons, then run npm run generate:books.
% canonical-source-hash: fnv1a64:648795848d5bf33f
% canonical-lessons: ML-C513-kyampu, ML-C513-duritasvasam, ML-C513-itiyuka, ML-C513-katapulakuka, ML-C513-valayuka

\chapter{After the Storm}
\label{ch:ml-after-the-storm}
\hlchaptermodality{513}

\begin{chapteropening}
\textbf{By the end of this chapter:} \emph{I can say a relief camp, disaster relief, to cave in, to be uprooted and to be in distress.}

Complete the last lesson of chapter 513: I can say a relief camp, disaster relief, to cave in, to be uprooted and to be in distress.
\end{chapteropening}

\section[\texorpdfstring{\ml{ക്യാമ്പ്} (kyāmpŭ)}{ക്യാമ്പ് (kyāmpŭ)}]{\ml{ക്യാമ്പ്} (kyāmpŭ) — a relief camp; a camp}

\label{lesson:ML-C513-kyampu}



\begin{quote}
Before the new one: say the Malayalam for to go under, then the Malayalam for to divert.
\end{quote}

\subsection*{You'll want to know: \ml{ക്യാമ്പ്}}
\textbf{\ml{ക്യാമ്പ്}} — \emph{kyāmpŭ} — \textquotedblleft{}a relief camp; a camp\textquotedblright{}.

\begin{quote}
\ml{വെള്ളപ്പൊക്കത്തിൽ} \ml{ആയിരം} \ml{പേർ} \ml{ക്യാമ്പിൽ} \ml{താമസിച്ചു}.
\end{quote}

\begin{quote}
\emph{veḷḷappokkattil āyiraṁ pēr kyāmpil tāmasiccu.}
\end{quote}

\begin{quote}
In the flood a thousand people stayed in the relief camp.
\end{quote}

Plural: \ml{ക്യാമ്പുകൾ}, \emph{kyāmpukaḷ}.

Schools often become \ml{ക്യാമ്പുകൾ} during a flood.

\begin{cousinweb}
\textbf{\ml{ക്യാമ്പ്}} (\emph{kyāmpŭ}) was borrowed from English \emph{camp}, a place of tents or shelters.
\end{cousinweb}

\begin{grammarlens}[title={What you've built}]
One more word for getting around.
\end{grammarlens}

\subsection*{Your turn}
\begin{itemize}
\raggedright
  \item \emph{Say it:} \emph{kyāmpŭ}
  \item \emph{Say it:} \emph{kyāmpŭ}, once more
  \item \emph{Say it:} \emph{veḷḷappokkattil āyiraṁ pēr kyāmpil tāmasiccu.} — In the flood a thousand people stayed in the relief camp.
  \item \emph{Recall:} say the Malayalam for to go under, then the Malayalam for to divert, then say \emph{kyāmpŭ} again
\end{itemize}

\begin{tcolorbox}[breakable,colback=teal!4,colframe=teal!35!black,title={Before you move on}]
What does \ml{ക്യാമ്പ്} mean? (\textquotedblleft{}A relief camp; a camp\textquotedblright{}.) Say it once more.
\end{tcolorbox}
\section[\texorpdfstring{\ml{ദുരിതാശ്വാസം} (duritāśvāsaṁ)}{ദുരിതാശ്വാസം (duritāśvāsaṁ)}]{\ml{ദുരിതാശ്വാസം} (duritāśvāsaṁ) — disaster relief}

\label{lesson:ML-C513-duritasvasam}



\begin{quote}
Before the new one: say the Malayalam for to divert, then the Malayalam for a relief camp.
\end{quote}

\subsection*{You'll want to know: \ml{ദുരിതാശ്വാസം}}
\textbf{\ml{ദുരിതാശ്വാസം}} — \emph{duritāśvāsaṁ} — \textquotedblleft{}disaster relief\textquotedblright{}.

\begin{quote}
\ml{നാട്ടുകാർ} \ml{ദുരിതാശ്വാസത്തിനായി} \ml{പണം} \ml{നൽകി}.
\end{quote}

\begin{quote}
\emph{nāṭṭukār duritāśvāsattināyi paṇaṁ nalki.}
\end{quote}

\begin{quote}
The local people gave money for disaster relief.
\end{quote}

\ml{ദുരിതം}, distress, plus \ml{ആശ്വാസം}, relief.

\begin{cousinweb}
\textbf{\ml{ദുരിതാശ്വാസം}} (\emph{duritāśvāsaṁ}) is a Malayalam compound built from two Sanskrit words, \emph{durita}, \textquotedblleft{}difficulty, distress\textquotedblright{} (Monier-Williams), and \emph{āśvāsa}, \textquotedblleft{}breathing freely\textquotedblright{} (CDIAL 1471).
\end{cousinweb}

\begin{grammarlens}[title={What you've built}]
One more word for getting around.
\end{grammarlens}

\subsection*{Your turn}
\begin{itemize}
\raggedright
  \item \emph{Say it:} \emph{duritāśvāsaṁ}
  \item \emph{Say it:} \emph{duritāśvāsaṁ}, once more
  \item \emph{Say it:} \emph{nāṭṭukār duritāśvāsattināyi paṇaṁ nalki.} — The local people gave money for disaster relief.
  \item \emph{Recall:} say the Malayalam for to divert, then the Malayalam for a relief camp, then say \emph{duritāśvāsaṁ} again
\end{itemize}

\begin{tcolorbox}[breakable,colback=teal!4,colframe=teal!35!black,title={Before you move on}]
What does \ml{ദുരിതാശ്വാസം} mean? (\textquotedblleft{}Disaster relief\textquotedblright{}.) Say it once more.
\end{tcolorbox}
\section[\texorpdfstring{\ml{ഇടിയുക} (iṭiyuka)}{ഇടിയുക (iṭiyuka)}]{\ml{ഇടിയുക} (iṭiyuka) — to cave in; to crumble}

\label{lesson:ML-C513-itiyuka}



\begin{quote}
Before the new one: say the Malayalam for a relief camp, then the Malayalam for disaster relief.
\end{quote}

\subsection*{You'll want to know: \ml{ഇടിയുക}}
\textbf{\ml{ഇടിയുക}} — \emph{iṭiyuka} — \textquotedblleft{}to cave in; to crumble\textquotedblright{}.

\begin{quote}
\ml{കനത്ത} \ml{മഴയിൽ} \ml{വഴിയുടെ} \ml{ഒരു} \ml{ഭാഗം} \ml{ഇടിഞ്ഞു}.
\end{quote}

\begin{quote}
\emph{kanatta maḻayil vaḻiyuṭe oru bhāgaṁ iṭiññu.}
\end{quote}

\begin{quote}
In the heavy rain part of the road caved in.
\end{quote}

Past: \ml{ഇടിഞ്ഞു}, \emph{iṭiññu}.

For earth, walls and banks: \ml{മണ്ണ്} \ml{ഇടിഞ്ഞു}, the earth gave way.

\begin{cousinweb}
\textbf{\ml{ഇടിയുക}} (\emph{iṭiyuka}) is native Dravidian, \textquotedblleft{}to crumble, to fall to pieces\textquotedblright{} (DEDR 432).
\end{cousinweb}

\begin{grammarlens}[title={What you've built}]
One more word for getting around.
\end{grammarlens}

\subsection*{Your turn}
\begin{itemize}
\raggedright
  \item \emph{Say it:} \emph{iṭiyuka}
  \item \emph{Say it:} \emph{iṭiyuka}, once more
  \item \emph{Say it:} \emph{kanatta maḻayil vaḻiyuṭe oru bhāgaṁ iṭiññu.} — In the heavy rain part of the road caved in.
  \item \emph{Recall:} say the Malayalam for a relief camp, then the Malayalam for disaster relief, then say \emph{iṭiyuka} again
\end{itemize}

\begin{tcolorbox}[breakable,colback=teal!4,colframe=teal!35!black,title={Before you move on}]
What does \ml{ഇടിയുക} mean? (\textquotedblleft{}To cave in; to crumble\textquotedblright{}.) Say it once more.
\end{tcolorbox}
\section[\texorpdfstring{\ml{കടപുഴകുക} (kaṭapuḻakuka)}{കടപുഴകുക (kaṭapuḻakuka)}]{\ml{കടപുഴകുക} (kaṭapuḻakuka) — to be uprooted}

\label{lesson:ML-C513-katapulakuka}



\begin{quote}
Before the new one: say the Malayalam for disaster relief, then the Malayalam for to cave in.
\end{quote}

\subsection*{You'll want to know: \ml{കടപുഴകുക}}
\textbf{\ml{കടപുഴകുക}} — \emph{kaṭapuḻakuka} — \textquotedblleft{}to be uprooted\textquotedblright{}.

\begin{quote}
\ml{കാറ്റിൽ} \ml{വഴിയോരത്തെ} \ml{ഒരു} \ml{വലിയ} \ml{മരം} \ml{കടപുഴകി}.
\end{quote}

\begin{quote}
\emph{kāṟṟil vaḻiyōratte oru valiya maraṁ kaṭapuḻaki.}
\end{quote}

\begin{quote}
In the wind a big tree by the roadside was uprooted.
\end{quote}

Past: \ml{കടപുഴകി}, \emph{kaṭapuḻaki}.

Only for trees torn out by the roots: \ml{മരം} \ml{കടപുഴകി}, the tree came down roots and all.

\begin{cousinweb}
\textbf{\ml{കടപുഴകുക}} (\emph{kaṭapuḻakuka}) is a Malayalam verb for a tree torn up by the roots; its origin is not settled.
\end{cousinweb}

\begin{grammarlens}[title={What you've built}]
One more word for getting around.
\end{grammarlens}

\subsection*{Your turn}
\begin{itemize}
\raggedright
  \item \emph{Say it:} \emph{kaṭapuḻakuka}
  \item \emph{Say it:} \emph{kaṭapuḻakuka}, once more
  \item \emph{Say it:} \emph{kāṟṟil vaḻiyōratte oru valiya maraṁ kaṭapuḻaki.} — In the wind a big tree by the roadside was uprooted.
  \item \emph{Recall:} say the Malayalam for disaster relief, then the Malayalam for to cave in, then say \emph{kaṭapuḻakuka} again
\end{itemize}

\begin{tcolorbox}[breakable,colback=teal!4,colframe=teal!35!black,title={Before you move on}]
What does \ml{കടപുഴകുക} mean? (\textquotedblleft{}To be uprooted\textquotedblright{}.) Say it once more.
\end{tcolorbox}
\section[\texorpdfstring{\ml{വലയുക} (valayuka)}{വലയുക (valayuka)}]{\ml{വലയുക} (valayuka) — to be in distress; to be put to hardship}

\label{lesson:ML-C513-valayuka}



\begin{quote}
Before the new one: say the Malayalam for to cave in, then the Malayalam for to be uprooted.
\end{quote}

\subsection*{You'll want to know: \ml{വലയുക}}
\textbf{\ml{വലയുക}} — \emph{valayuka} — \textquotedblleft{}to be in distress; to be put to hardship\textquotedblright{}.

\begin{quote}
\ml{ബസ്} \ml{ഇല്ലാതെ} \ml{യാത്രക്കാർ} \ml{വലഞ്ഞു}.
\end{quote}

\begin{quote}
\emph{basŭ illāte yātrakkār valaññu.}
\end{quote}

\begin{quote}
With no bus, the passengers were stranded and suffered.
\end{quote}

Past: \ml{വലഞ്ഞു}, \emph{valaññu}.

Common in the news: \ml{യാത്രക്കാർ} \ml{വലഞ്ഞു}, the passengers had a hard time.

\begin{cousinweb}
\textbf{\ml{വലയുക}} (\emph{valayuka}) is native Dravidian, \textquotedblleft{}to be straitened, pressed, distressed\textquotedblright{} (DEDR 5281).
\end{cousinweb}

\begin{grammarlens}[title={What you've built}]
One more word for getting around.
\end{grammarlens}

\subsection*{Your turn}
\begin{itemize}
\raggedright
  \item \emph{Say it:} \emph{valayuka}
  \item \emph{Say it:} \emph{valayuka}, once more
  \item \emph{Say it:} \emph{basŭ illāte yātrakkār valaññu.} — With no bus, the passengers were stranded and suffered.
  \item \emph{Recall:} say the Malayalam for to cave in, then the Malayalam for to be uprooted, then say \emph{valayuka} again
\end{itemize}

\begin{tcolorbox}[breakable,colback=teal!4,colframe=teal!35!black,title={Before you move on}]
What does \ml{വലയുക} mean? (\textquotedblleft{}To be in distress; to be put to hardship\textquotedblright{}.) Say it once more.
\end{tcolorbox}
